\section*{Chapitre \ref{ch:ll:managed-and-functional-languages-on-the-desk} --- Langages gérés et fonctionnels sur le desk}

\begin{solution}{exo:ll:managed-and-functional-languages-on-the-desk:1}
$64 \times 100\,000 = 6.4$~Mo par seconde remplit 1~GiB en 167.8 secondes~: environ 139 collections de jeune génération sur une session de 6.5~heures
(voir la \cref{prop:ll:managed-and-functional-languages-on-the-desk:budget}).
\end{solution}

\begin{solution}{exo:ll:managed-and-functional-languages-on-the-desk:2}
$0.001 \times 200\,000 = 200$ messages~; le dernier arrive juste avant la fin de la pause et le premier attend environ une milliseconde,
et tous sont ensuite traités en retard à leur tour, si bien que la file met encore 200 durées de traitement pour se vider.
\end{solution}

\begin{solution}{exo:ll:managed-and-functional-languages-on-the-desk:3}
Sur la mesure retenue, le maximum passe d'environ \qty{19.6}{\milli\second} à \qty{0.58}{\milli\second}, soit
environ trente fois, et le p99.9 passe d'environ \qty{26}{\micro\second} à \qty{1.7}{\micro\second}~; la médiane se dégrade légèrement,
\qty{335}{\nano\second} contre \qty{313}{\nano\second}.
\end{solution}

\begin{solution}{exo:ll:managed-and-functional-languages-on-the-desk:4}
CPython déclenche une collection de jeune génération lorsque l'écart entre les allocations et les désallocations d'objets suivis dépasse le seuil (700 par
défaut sous Python 3.10). Une deque bornée libère un ancien ordre pour chaque nouvel ordre, si bien que l'écart ne croît jamais~; un carnet qui grossit
ajoute plus qu'il ne libère et franchit le seuil tous les quelques centaines de messages.
\end{solution}

\begin{solution}{exo:ll:managed-and-functional-languages-on-the-desk:5}
Le chemin critique est interprété ou peu compilé tant que le compilateur à la volée (JIT) n'a pas vu assez d'exécutions~; l'ouverture est
aussi le moment où des chemins rares sont empruntés pour la première fois, et peuvent être désoptimisés~. Remèdes~: chauffer le code en rejouant le flux enregistré à travers le vrai code
face à un simulateur avant l'ouverture~; garder stables les types et les branchements du chemin critique~; ou compiler à l'avance.
\end{solution}

\begin{solution}{exo:ll:managed-and-functional-languages-on-the-desk:6}
$20 \times 10^6 / 6 \times 10^6 \approx 3.33$ secondes de flux entrant à plein débit.
\end{solution}

\begin{solution}{exo:ll:managed-and-functional-languages-on-the-desk:7}
Le gel déplace chaque objet vivant à cet instant vers une génération permanente que les collections ignorent, si bien que les collections complètes
ne parcourent que les ordres ajoutés depuis le dernier gel et sont bien plus courtes. Le prix se paie en mémoire~: les objets gelés qui meurent
ne sont jamais récupérés tant que le tas n'est pas dégelé, si bien qu'un carnet dont les ordres sont annulés conserve leur mémoire.
\end{solution}

\begin{solution}{exo:ll:managed-and-functional-languages-on-the-desk:8}
Moins d'une milliseconde représente mille microsecondes, bien au-delà d'un budget d'une microseconde~; et un \omterm{def:ll:managed-and-functional-languages-on-the-desk:gc}{ramasse-miettes} concurrent prélève encore
des cœurs, de la bande passante mémoire et du cache au programme pendant son exécution, et suspend les threads aux safepoints. Il affecte la latence
sauf si la passerelle n'alloue rien pendant la session.
\end{solution}

\begin{solution}{pb:ll:managed-and-functional-languages-on-the-desk:1}
\begin{enumerate}
\item $48 \times 200\,000 = 9.6$~Mo par seconde~; 224.64~Go sur une session.
\item Toutes les $2~\text{GiB}/9.6~\text{Mo/s} = 223.7$ secondes~: environ 105 fois par session.
\item $2~\text{GiB}/(200\,000 \times 23\,400) = 0.459$ octets.
\item Zéro~: aucune allocation par message sur le chemin critique.
\item $105 \times \qty{0.5}{\milli\second} \approx \qty{52}{\milli\second}$.
\item $0.0005 \times 200\,000 = 100$ messages~; le premier attend environ \qty{0.5}{\milli\second}, et la file met
encore 100 durées de traitement pour se vider.
\item Le taux d'allocation triple~: environ 314 collections et \qty{157}{\milli\second} de pause, avec 300 messages en attente derrière chaque
pause.
\item Les files et les pauses croissent avec le trafic, et les journées les plus chargées sont celles où la latence compte le plus.
\item Des champs primitifs et des collections primitives (tableaux de \texttt{long}), sans boxing.
\item Un journal binaire~: le chemin critique écrit un identifiant de format et les arguments bruts dans un tampon circulaire préalloué~; la mise en forme se fait
ailleurs, plus tard.
\item En \omterm{def:ll:managed-and-functional-languages-on-the-desk:offheap}{mémoire hors tas}, ou dans des tableaux plats de primitives alloués au démarrage, de sorte qu'aucune collection ne les parcoure ni ne les déplace.
\item Mesurer les octets alloués par thread avant et après un rejeu de messages (les compteurs d'allocation par thread de la JVM, ou
un profileur d'allocation) et vérifier qu'il est nul après la \omterm{def:ll:managed-and-functional-languages-on-the-desk:jit}{période de chauffe}.
\item \textbf{Résultat nommé.} Le moteur peut allouer au plus 0.459 octets par message, en pratique rien, pour ne collecter qu'une fois par
session~; à 48 octets par message, il collecte environ 105 fois et passe environ \qty{52}{\milli\second} en pause.
\item Les safepoints et les threads d'arrière-plan du \omterm{def:ll:managed-and-functional-languages-on-the-desk:gc}{ramasse-miettes}, qui prennent encore des cœurs et du cache~; et le coût
du démarrage et de la \omterm{def:ll:managed-and-functional-languages-on-the-desk:jit}{période de chauffe}.
\item Parce que le compilateur à la volée (JIT) doit encore voir le chemin critique exécuté avant de l'optimiser.
\item Que supprimer l'allocation par message retire le \omterm{def:ll:managed-and-functional-languages-on-the-desk:gc}{ramasse-miettes} de la boucle, et que le pire message est la
collection complète sur le tas vivant.
\item Du code lisible, correct et efficace, écrit et modifié vite~: la correction par ingénieur compte davantage que la vitesse par message.
\item Lorsque le budget se situe à quelques microsecondes à un percentile élevé, et que l'équipe ne peut pas garantir une allocation nulle dans
chaque version.
\item Le \omterm{def:ll:where-latency-comes-from:budget}{budget de latence} au percentile qui compte, l'allocation par message et les pauses mesurées pendant une
journée chargée rejouée.
\item Une pause proportionnelle à ce qu'alloue le chemin critique et à ce que le tas garde vivant.
\end{enumerate}
\end{solution}

\begin{solution}{iq:ll:managed-and-functional-languages-on-the-desk:1}
Il parcourt les objets atteignables depuis les racines et récupère le reste. La plupart des objets meurent jeunes, si bien que les nouveaux objets vont dans une petite génération jeune
collectée souvent, à faible coût, en proportion des survivants~; les survivants passent dans une génération âgée collectée
rarement.
\iqlookfor{l'hypothèse générationnelle faible et le coût proportionnel aux survivants.}
\end{solution}

\begin{solution}{iq:ll:managed-and-functional-languages-on-the-desk:2}
Ne rien allouer par message après la \omterm{def:ll:managed-and-functional-languages-on-the-desk:jit}{période de chauffe}~: pools préalloués et tampons circulaires, tableaux de primitives au lieu de
collections boxées~; \omterm{def:ll:managed-and-functional-languages-on-the-desk:offheap}{mémoire hors tas} pour l'état volumineux, journal binaire~; un collecteur discret et une génération jeune dimensionnée pour le
résiduel~; \omterm{def:ll:managed-and-functional-languages-on-the-desk:jit}{période de chauffe} avant l'ouverture~; mesurer l'allocation dans les tests.
\iqlookfor{aucune allocation, état hors tas, \omterm{def:ll:managed-and-functional-languages-on-the-desk:jit}{période de chauffe}, mesure.}
\end{solution}

\begin{solution}{iq:ll:managed-and-functional-languages-on-the-desk:3}
Le temps qui s'écoule avant que le JIT ait compilé le chemin critique avec son profil~; avant l'ouverture, rejouer le flux enregistré à travers le
vrai code (ordres envoyés à un simulateur) jusqu'à ce que les méthodes chaudes soient compilées, et garder le profil représentatif.
\iqlookfor{\omterm{def:ll:managed-and-functional-languages-on-the-desk:jit}{période de chauffe} avec un trafic réaliste, jamais avec de vrais ordres.}
\end{solution}

\begin{solution}{iq:ll:managed-and-functional-languages-on-the-desk:4}
Le l'environnement d'exécution rejette le code optimisé dont les hypothèses (un type unique à un site d'appel, une branche jamais prise) ont été invalidées,
et se replie sur un code plus lent~; le premier événement rare de la journée, souvent dans un moment de forte volatilité, est ce qui les invalide.
\iqlookfor{optimisation spéculative et son repli.}
\end{solution}

\begin{solution}{iq:ll:managed-and-functional-languages-on-the-desk:5}
Correction et vitesse de développement~: des types expressifs qui encodent des invariants (chaque cas d'un état d'ordre traité), des données
immuables par défaut, un code concis que les traders peuvent relire, comme l'a soutenu la société qui a publié son retour d'expérience en 2008~; avec
des performances suffisantes pour son budget.
\iqlookfor{correction et relecture, appuyées par le retour d'expérience publié.}
\end{solution}

\begin{solution}{iq:ll:managed-and-functional-languages-on-the-desk:6}
Enregistrer la durée de chaque collection avec les callbacks du module \texttt{gc} et les corréler aux messages lents~; compter
les collections par génération. Puis arrêter d'allouer par message (tableaux préalloués), geler les objets à longue durée de vie après le démarrage,
ou programmer explicitement des collections complètes dans les moments calmes.
\iqlookfor{mesurer le \omterm{def:ll:managed-and-functional-languages-on-the-desk:gc}{ramasse-miettes}, puis supprimer sa cause.}
\end{solution}
